% Phần 2: Mô hình v2, cách đánh giá, kết quả chung, Yêu cầu 1.
% Mọi con số lấy qua \so{...} (so_lieu.tex, sinh tự động). Không gõ số tay.

\section{Mô hình v2}
\label{sec:p2-mo-hinh}

\subsection{Vì sao chọn encoder thay cho mô hình sinh}

Kế hoạch ban đầu là fine-tune một mô hình sinh nhỏ (Qwen3-1.7B, QLoRA) để viết ra JSON. Khi làm thật, nhóm huấn luyện
song song 4 ứng viên trên cùng dữ liệu và cùng seed, rồi chọn trên tập dev (so sánh đầy đủ ở phần Yêu cầu 3). Ứng viên
dựa trên LLM 0,6B (GuardMerge) thua encoder có ý nghĩa thống kê (kiểm định McNemar, $p = \so{dev.mc.guard}$) và chậm
gấp hơn hai lần (\so{dev.guard.p50} ms so với \so{dev.mmbert.p50} ms mỗi câu trên GPU). Vì vậy v2 dùng encoder.

\subsection{Cấu trúc}

Mô hình chốt có tên nội bộ \textbf{SafeVi-DualORC-CRF}, xây trên \textbf{mmBERT-base} (encoder đa ngôn ngữ). Nó đọc
cả đoạn hội thoại một lần (tối đa \so{v2.maxtok} token), không chỉ dòng cuối, và có hai đầu ra:

\begin{enumerate}
  \item \textbf{Đầu phân loại:} trả xác suất cho 5 nhãn. Nhãn có xác suất cao nhất là kết quả.
  \item \textbf{Đầu gắn thẻ:} gán cho từng token một thẻ B (bắt đầu cụm), I (trong cụm) hoặc O (ngoài cụm), dùng thêm
        tầng CRF để các thẻ nối với nhau hợp lệ. Các đoạn B--I liền nhau là cụm bằng chứng.
\end{enumerate}

Vì cụm bằng chứng được \emph{cắt ra từ chính văn bản}, nó luôn khớp nguyên văn. Ngoài encoder, mô hình có thêm một nhánh CNN đọc từng byte của văn bản, nhằm nhận ra chữ
bị thay bằng ký tự lạ mà bộ tách từ không quen. Hai đầu ra được ghép thành JSON \texttt{\{"nhan": ..., "bang\_chung":
[...]\}} đúng định dạng đã nêu ở phần nhãn.

\begin{figure}[H]
\centering
\begin{tikzpicture}[
  font=\small,
  hop/.style={draw=navy, rounded corners=2pt, fill=nenxanh, align=center, text width=2.5cm, minimum height=1.1cm},
  tuychon/.style={hop, dashed, fill=nenxam},
  mui/.style={-{Stealth[length=2mm]}, thick, navy}
]
  \node[hop] (vao) {Đoạn hội thoại\\(1--8 dòng)};
  \node[tuychon, right=0.5cm of vao] (chuan) {Chuẩn hóa\\(tùy chọn)};
  \node[hop, right=0.5cm of chuan, text width=2.7cm] (enc) {mmBERT-base\\+ CNN đọc byte};
  \node[hop, right=0.7cm of enc, yshift=0.8cm, text width=3.0cm] (pl) {Đầu phân loại\\xác suất 5 nhãn};
  \node[hop, right=0.7cm of enc, yshift=-0.8cm, text width=3.0cm] (gt) {Đầu gắn thẻ\\BIO + CRF $\to$ cụm};
  % Hộp JSON đặt hẳn dưới đầu gắn thẻ, mép phải thẳng hàng, để không chồng lên hộp phía trên
  \node[hop, anchor=north east, text width=4.6cm] (ra) at ([yshift=-0.9cm]gt.south east)
    {JSON: \texttt{nhan} + \texttt{bang\_chung}};
  \draw[mui] (vao) -- (chuan);
  \draw[mui] (chuan) -- (enc);
  \draw[mui] (enc.east) -- ++(0.3,0) |- (pl.west);
  \draw[mui] (enc.east) -- ++(0.3,0) |- (gt.west);
  \draw[mui] (pl.east) -- ++(0.4,0) |- (ra.east);
  \draw[mui] (gt.east) -- ++(0.4,0) |- (ra.east);
\end{tikzpicture}
\caption{Luồng xử lý của v2. Hai đầu ra dùng chung encoder nhưng có ngưỡng và cách chấm riêng.}
\label{fig:p2-luong}
\end{figure}

\subsection{Huấn luyện}

Tệp train lưu bản gốc C0. Khi huấn luyện, mỗi câu được biến dạng ngẫu nhiên ngay lúc đọc: bỏ dấu với xác suất
\so{v2.accent}, biến dạng bề mặt thường gặp (kiểu C1 đã thấy) với xác suất \so{v2.surface}. Mô hình còn được phạt nếu dự
đoán trên bản gốc và bản biến dạng khác nhau, và được huấn luyện đối kháng FGM như v1 (cộng nhiễu nhỏ vào embedding,
$\varepsilon = \so{v2.fgm}$). Các kiểu điểm mù C2 và nhóm C1 giữ ngoài train không được dùng ở bước này.
Huấn luyện với \so{v2.epoch} epoch, seed 42.

\subsection{Chia dữ liệu và các ngưỡng}

Tập train được chia thành ba phần: \textbf{fit} (\so{dev.chia.fit} mẫu) để học, \textbf{dev} (\so{dev.chia.dev} mẫu) để
chọn mô hình và checkpoint, \textbf{cal} (\so{dev.chia.cal} mẫu) để hiệu chỉnh. Dev và cal tách từ train nên nhãn của
chúng cũng do LLM gán; không câu test nào được dùng ở đây. Trên cal, nhóm dò và khóa ba tham số:

\begin{itemize}
  \item nhiệt độ $T = \so{v2.T}$ để xác suất phân loại sát với tỉ lệ đúng thật (chia logit cho $T$ trước softmax);
  \item ngưỡng bằng chứng \so{v2.nguongbc}: token có xác suất thuộc cụm từ ngưỡng này trở lên mới được tô;
  \item ngưỡng riêng \so{v2.nguongmottu} cho cụm chỉ có một từ, cao hơn để tránh tô lẻ tẻ.
\end{itemize}

Cả ba được ghi vào \tep{calibration.json} và đóng băng trước khi chạy test.

% ---------------------------------------------------------------------------
\section{Cách đánh giá và kết quả chung}
\label{sec:p2-danh-gia}

\subsection{Ba hệ được so sánh}

\begin{itemize}
  \item \textbf{v1:} SafeViChi vòng loại (ViHateT5, HATE/CLEAN). Chỉ có nhãn nhị phân, không có bằng chứng, nên chỉ chấm
        được phần nguy hại/an toàn. HATE được tính là ``nguy hại'', CLEAN là \nhan{AN_TOAN}.
  \item \textbf{gpt-oss-120b qua API:} LLM lớn, zero-shot. Prompt chính là hướng dẫn gán nhãn của nhóm (định nghĩa 5 nhãn,
        thứ tự quyết định, ví dụ, quy tắc bằng chứng). Mức suy luận \texttt{\so{api.reasoning}}.
  \item \textbf{v2:} mô hình ở Mục~\ref{sec:p2-mo-hinh}, không bật chuẩn hóa đầu vào.
\end{itemize}

Cả ba chạy trên cùng bộ test: C0 (\so{C0.v2.n} đoạn gốc do người viết), C1 (\so{C1.v2.n} bản biến dạng thường gặp như
bỏ dấu, teencode, gõ sai) và C2 (\so{C2.v2.n} bản biến dạng ``điểm mù'' chưa từng xuất hiện khi train v2).

\subsection{Độ đo}

\begin{itemize}
  \item \textbf{Macro-F1:} trung bình F1 của 5 nhãn, mỗi nhãn quan trọng như nhau. Đo khả năng gán nhãn đúng tình huống
  \item \textbf{Binary-F1:} gộp 4 nhãn nguy hại thành một, đo khả năng phân biệt nguy hại với an toàn.
  \item \textbf{FPR:} tỉ lệ câu \nhan{AN_TOAN} bị báo nguy hại (báo nhầm). Càng thấp càng tốt.
  \item \textbf{BC P/R/F1 theo ý:} chấm bằng chứng. Mỗi cụm người gán là một ý, precision là phần cụm mô hình tô
        trúng ý đúng, recall là phần ý đúng được tô ra.
\end{itemize}

Trong toàn báo cáo, F1, precision và recall được nhân 100 (ví dụ 86,5); các tỉ lệ ghi kèm dấu \%. Khoảng tin cậy 95\%
tính bằng bootstrap 1.000 lần, seed 42.

\subsection{Kết quả}

\begin{table}[H]
\centering
\small
\caption{Kết quả chính trên bộ test. ``--'': hệ không có đầu ra tương ứng.}
\label{tab:p2-chinh}
\begin{tabular}{llccccc}
\toprule
Tập & Hệ & Macro-F1 & KTC 95\% Macro-F1 & Binary-F1 & FPR & BC F1 theo ý \\
\midrule
\multirow{3}{*}{C0} & v1 & -- & -- & \so{C0.v1.bin} & \so{C0.v1.fpr} & -- \\
 & gpt-oss-120b & \so{C0.llm.macro} & \so{C0.llm.macro.lo} -- \so{C0.llm.macro.hi} & \so{C0.llm.bin} & \so{C0.llm.fpr} & \so{C0.llm.bcf} \\
 & \textbf{v2} & \textbf{\so{C0.v2.macro}} & \so{C0.v2.macro.lo} -- \so{C0.v2.macro.hi} & \textbf{\so{C0.v2.bin}} & \so{C0.v2.fpr} & \textbf{\so{C0.v2.bcf}} \\
\midrule
\multirow{3}{*}{C1} & v1 & -- & -- & \so{C1.v1.bin} & \so{C1.v1.fpr} & -- \\
 & gpt-oss-120b & \so{C1.llm.macro} & \so{C1.llm.macro.lo} -- \so{C1.llm.macro.hi} & \so{C1.llm.bin} & \so{C1.llm.fpr} & \so{C1.llm.bcf} \\
 & \textbf{v2} & \textbf{\so{C1.v2.macro}} & \so{C1.v2.macro.lo} -- \so{C1.v2.macro.hi} & \textbf{\so{C1.v2.bin}} & \so{C1.v2.fpr} & \textbf{\so{C1.v2.bcf}} \\
\midrule
\multirow{3}{*}{C2} & v1 & -- & -- & \so{C2.v1.bin} & \so{C2.v1.fpr} & -- \\
 & gpt-oss-120b & \so{C2.llm.macro} & \so{C2.llm.macro.lo} -- \so{C2.llm.macro.hi} & \so{C2.llm.bin} & \so{C2.llm.fpr} & \so{C2.llm.bcf} \\
 & \textbf{v2} & \textbf{\so{C2.v2.macro}} & \so{C2.v2.macro.lo} -- \so{C2.v2.macro.hi} & \textbf{\so{C2.v2.bin}} & \so{C2.v2.fpr} & \textbf{\so{C2.v2.bcf}} \\
\bottomrule
\end{tabular}
\end{table}

\begin{figure}[H]
\centering
% Sinh tự động bằng `python -m src.v2.eval.xuat_so_lieu_bao_cao`. Không sửa tay.
\begin{tikzpicture}
\begin{axis}[ybar=3pt, bar width=18pt, width=\linewidth, height=6cm, ymin=0, ymax=100,
  symbolic x coords={C0,C1,C2}, xtick=data, enlarge x limits=0.25, ylabel={Binary-F1},
  axis x line*=bottom, axis y line*=left, ymajorgrids, grid style={draw=black!12},
  tick label style={font=\small}, label style={font=\small},
  nodes near coords={\pgfmathprintnumber[fixed,fixed zerofill,precision=1,use comma]{\pgfplotspointmeta}},
  every node near coord/.append style={font=\scriptsize, color=black!75},
  legend style={at={(0.5,1.03)}, anchor=south, legend columns=3, draw=none, font=\small,
    /tikz/every even column/.append style={column sep=0.4cm}}]
\addplot[fill=mauvmot, draw=white, area legend] coordinates {(C0,31.2) (C1,31.0) (C2,39.5)};
\addlegendentry{SafeViChi v1}
\addplot[fill=maullm, draw=white, area legend] coordinates {(C0,89.2) (C1,74.5) (C2,77.9)};
\addlegendentry{gpt-oss-120b (API)}
\addplot[fill=mauvhai, draw=white, area legend] coordinates {(C0,92.9) (C1,90.3) (C2,90.9)};
\addlegendentry{SafeViChi v2}
\end{axis}
\end{tikzpicture}

\caption{Binary-F1 (nguy hại/an toàn) của ba hệ trên C0, C1, C2. Số đầy đủ ở Bảng~\ref{tab:p2-chinh}.}
\label{fig:p2-binary}
\end{figure}

\textbf{Nhận xét.}
\begin{itemize}
  \item \textbf{v2 tốt nhất trên cả ba tập.} Khoảng tin cậy 95\% của Macro-F1 giữa v2 và gpt-oss-120b không chồng nhau ở
        cả C0, C1 và C2, nên khoảng cách không phải do may rủi của mẫu.
  \item \textbf{v2 giữ điểm khi chữ bị biến dạng.} Binary-F1 của v2 đi từ \so{C0.v2.bin} (C0) xuống \so{C2.v2.bin} (C2);
        của gpt-oss-120b từ \so{C0.llm.bin} xuống \so{C2.llm.bin}.
  % \item \textbf{v1 gần như chỉ bắt được xúc phạm.} Ở C0, v1 báo HATE cho \so{C0.v1.xp.r} câu \nhan{XUC_PHAM}, nhưng chỉ
  %       \so{C0.v1.dd.r} câu \nhan{DE_DOA}, \so{C0.v1.qr.r} câu \nhan{QUAY_ROI} và \so{C0.v1.ld.r} câu \nhan{LUA_DAO}.
  %       Precision \so{C0.v1.p} của v1 chỉ phản ánh việc nó hiếm khi báo. Đây là lý do phải đổi bộ nhãn.
  \item \textbf{Đổi lại, v2 báo nhầm nhiều hơn.} Ở C0, FPR của v2 là \so{C0.v2.fpr}, của gpt-oss-120b là
        \so{C0.llm.fpr}. gpt-oss-120b thận trọng hơn: precision nguy hại \so{C0.llm.p} nhưng recall chỉ \so{C0.llm.r}
        (v2: \so{C0.v2.p} và \so{C0.v2.r}).
\end{itemize}

% ---------------------------------------------------------------------------
\section{Yêu cầu 1: giảm kiểm duyệt sai và giải thích sai}
\label{sec:p2-yc1}

\subsection{Tách kết quả phân loại khỏi bằng chứng}

Ở v1, bằng chứng tìm bằng occlusion \emph{sau} khi đã phân loại: che từng cụm rồi xem điểm giảm bao nhiêu. Cách này phải
chạy lại mô hình nhiều lần và chất lượng bằng chứng gắn chặt với bộ phân loại. v2 tách hẳn hai việc:

\begin{itemize}
  \item \textbf{Hai đầu ra riêng}, mỗi đầu có hàm mất mát riêng khi train.
  \item \textbf{Hai ngưỡng riêng:} nhiệt độ $T$ cho xác suất nhãn; ngưỡng \so{v2.nguongbc} và \so{v2.nguongmottu} cho
        cụm bằng chứng. Đổi ngưỡng bằng chứng không làm đổi nhãn.
  \item \textbf{Hai bộ độ đo riêng:} nhãn chấm bằng P/R/F1 theo nhãn; bằng chứng chấm bằng P/R/F1 theo ý.
  \item \textbf{Vai trò khác nhau trong ứng dụng:} nhãn quyết định ứng dụng làm gì (làm mờ, cảnh báo, gợi ý chặn); bằng
        chứng chỉ để người dùng hiểu vì sao, nên chỉ có nghĩa khi nhãn là nguy hại.
\end{itemize}

\subsection{Kết quả theo từng nhãn}

\begin{table}[H]
\centering
\small
\caption{Kết quả theo nhãn trên C0. n: số câu có nhãn đúng là nhãn đó.}
\label{tab:p2-theo-nhan}
\begin{tabular}{lcccccccc}
\toprule
 & & \multicolumn{4}{c}{v2} & \multicolumn{2}{c}{gpt-oss-120b} \\
\cmidrule(lr){3-6}\cmidrule(lr){7-8}
Nhãn & n & P & R & F1 & BC F1 & F1 & BC F1 \\
\midrule
\nhan{AN_TOAN}  & \so{C0.v2.at.n} & \so{C0.v2.at.p} & \so{C0.v2.at.r} & \so{C0.v2.at.f} & -- & \so{C0.llm.at.f} & -- \\
\nhan{XUC_PHAM} & \so{C0.v2.xp.n} & \so{C0.v2.xp.p} & \so{C0.v2.xp.r} & \so{C0.v2.xp.f} & \so{C0.v2.xp.bcf} & \so{C0.llm.xp.f} & \so{C0.llm.xp.bcf} \\
\nhan{QUAY_ROI} & \so{C0.v2.qr.n} & \so{C0.v2.qr.p} & \so{C0.v2.qr.r} & \so{C0.v2.qr.f} & \so{C0.v2.qr.bcf} & \so{C0.llm.qr.f} & \so{C0.llm.qr.bcf} \\
\nhan{DE_DOA}   & \so{C0.v2.dd.n} & \so{C0.v2.dd.p} & \so{C0.v2.dd.r} & \so{C0.v2.dd.f} & \so{C0.v2.dd.bcf} & \so{C0.llm.dd.f} & \so{C0.llm.dd.bcf} \\
\nhan{LUA_DAO}  & \so{C0.v2.ld.n} & \so{C0.v2.ld.p} & \so{C0.v2.ld.r} & \so{C0.v2.ld.f} & \so{C0.v2.ld.bcf} & \so{C0.llm.ld.f} & \so{C0.llm.ld.bcf} \\
\bottomrule
\end{tabular}
\end{table}

\nhan{LUA_DAO} và \nhan{XUC_PHAM} là hai nhãn tốt nhất. \nhan{QUAY_ROI} yếu nhất, chủ yếu do recall thấp: mô hình bỏ sót
khá nhiều câu quấy rối. So với gpt-oss-120b, v2 hơn rõ ở \nhan{XUC_PHAM}, \nhan{QUAY_ROI}, \nhan{DE_DOA} và ngang ở
\nhan{LUA_DAO}; bằng chứng của v2 tốt hơn ở cả bốn nhãn nguy hại.

\begin{figure}[H]
\centering
% Sinh tự động bằng `python -m src.v2.eval.xuat_so_lieu_bao_cao`. Không sửa tay.
\begin{tikzpicture}
\fill[mauvmot!90] (0.0,0.0) rectangle ++(1.45,-1.45);\node[font=\small\bfseries,text=white] at (0.725,-0.725) {86};
\fill[mauvmot!0] (1.45,0.0) rectangle ++(1.45,-1.45);\node[font=\small\bfseries,text=black] at (2.175,-0.725) {0};
\fill[mauvmot!1] (2.9,0.0) rectangle ++(1.45,-1.45);\node[font=\small\bfseries,text=black] at (3.625,-0.725) {1};
\fill[mauvmot!6] (4.35,0.0) rectangle ++(1.45,-1.45);\node[font=\small\bfseries,text=black] at (5.074999999999999,-0.725) {6};
\fill[mauvmot!3] (5.8,0.0) rectangle ++(1.45,-1.45);\node[font=\small\bfseries,text=black] at (6.5249999999999995,-0.725) {3};
\node[anchor=east,font=\footnotesize] at (-0.1,-0.725) {\nhan{AN_TOAN}};
\fill[mauvmot!4] (0.0,-1.45) rectangle ++(1.45,-1.45);\node[font=\small\bfseries,text=black] at (0.725,-2.175) {2};
\fill[mauvmot!91] (1.45,-1.45) rectangle ++(1.45,-1.45);\node[font=\small\bfseries,text=white] at (2.175,-2.175) {48};
\fill[mauvmot!6] (2.9,-1.45) rectangle ++(1.45,-1.45);\node[font=\small\bfseries,text=black] at (3.625,-2.175) {3};
\fill[mauvmot!0] (4.35,-1.45) rectangle ++(1.45,-1.45);\node[font=\small\bfseries,text=black] at (5.074999999999999,-2.175) {0};
\fill[mauvmot!0] (5.8,-1.45) rectangle ++(1.45,-1.45);\node[font=\small\bfseries,text=black] at (6.5249999999999995,-2.175) {0};
\node[anchor=east,font=\footnotesize] at (-0.1,-2.175) {\nhan{XUC_PHAM}};
\fill[mauvmot!16] (0.0,-2.9) rectangle ++(1.45,-1.45);\node[font=\small\bfseries,text=black] at (0.725,-3.625) {8};
\fill[mauvmot!0] (1.45,-2.9) rectangle ++(1.45,-1.45);\node[font=\small\bfseries,text=black] at (2.175,-3.625) {0};
\fill[mauvmot!71] (2.9,-2.9) rectangle ++(1.45,-1.45);\node[font=\small\bfseries,text=white] at (3.625,-3.625) {36};
\fill[mauvmot!14] (4.35,-2.9) rectangle ++(1.45,-1.45);\node[font=\small\bfseries,text=black] at (5.074999999999999,-3.625) {7};
\fill[mauvmot!0] (5.8,-2.9) rectangle ++(1.45,-1.45);\node[font=\small\bfseries,text=black] at (6.5249999999999995,-3.625) {0};
\node[anchor=east,font=\footnotesize] at (-0.1,-3.625) {\nhan{QUAY_ROI}};
\fill[mauvmot!15] (0.0,-4.35) rectangle ++(1.45,-1.45);\node[font=\small\bfseries,text=black] at (0.725,-5.074999999999999) {9};
\fill[mauvmot!2] (1.45,-4.35) rectangle ++(1.45,-1.45);\node[font=\small\bfseries,text=black] at (2.175,-5.074999999999999) {1};
\fill[mauvmot!2] (2.9,-4.35) rectangle ++(1.45,-1.45);\node[font=\small\bfseries,text=black] at (3.625,-5.074999999999999) {1};
\fill[mauvmot!82] (4.35,-4.35) rectangle ++(1.45,-1.45);\node[font=\small\bfseries,text=white] at (5.074999999999999,-5.074999999999999) {51};
\fill[mauvmot!0] (5.8,-4.35) rectangle ++(1.45,-1.45);\node[font=\small\bfseries,text=black] at (6.5249999999999995,-5.074999999999999) {0};
\node[anchor=east,font=\footnotesize] at (-0.1,-5.074999999999999) {\nhan{DE_DOA}};
\fill[mauvmot!6] (0.0,-5.8) rectangle ++(1.45,-1.45);\node[font=\small\bfseries,text=black] at (0.725,-6.5249999999999995) {4};
\fill[mauvmot!0] (1.45,-5.8) rectangle ++(1.45,-1.45);\node[font=\small\bfseries,text=black] at (2.175,-6.5249999999999995) {0};
\fill[mauvmot!0] (2.9,-5.8) rectangle ++(1.45,-1.45);\node[font=\small\bfseries,text=black] at (3.625,-6.5249999999999995) {0};
\fill[mauvmot!0] (4.35,-5.8) rectangle ++(1.45,-1.45);\node[font=\small\bfseries,text=black] at (5.074999999999999,-6.5249999999999995) {0};
\fill[mauvmot!94] (5.8,-5.8) rectangle ++(1.45,-1.45);\node[font=\small\bfseries,text=white] at (6.5249999999999995,-6.5249999999999995) {68};
\node[anchor=east,font=\footnotesize] at (-0.1,-6.5249999999999995) {\nhan{LUA_DAO}};
\node[anchor=west,rotate=40,font=\footnotesize] at (0.5249999999999999,0.15) {\nhan{AN_TOAN}};
\node[anchor=west,rotate=40,font=\footnotesize] at (1.9749999999999999,0.15) {\nhan{XUC_PHAM}};
\node[anchor=west,rotate=40,font=\footnotesize] at (3.425,0.15) {\nhan{QUAY_ROI}};
\node[anchor=west,rotate=40,font=\footnotesize] at (4.874999999999999,0.15) {\nhan{DE_DOA}};
\node[anchor=west,rotate=40,font=\footnotesize] at (6.324999999999999,0.15) {\nhan{LUA_DAO}};
\draw[black!25] (0,0) grid[step=1.45] (7.25,-7.25);
\node[font=\small,rotate=90] at (-3.0,-3.625) {Nhãn đúng};
\node[font=\small] at (3.625,2.0) {Nhãn SafeViChi v2 dự đoán};
\end{tikzpicture}

\caption{Ma trận nhầm lẫn của v2 trên C0. Số trong ô là số câu; màu đậm theo tỉ lệ trong hàng. Ngoài đường chéo,
lỗi dồn vào cột \nhan{AN_TOAN} (bỏ sót) và cột \nhan{DE_DOA} (báo nhầm, lẫn quấy rối với đe dọa).}
\label{fig:p2-nham-lan}
\end{figure}

Ngoài nhầm với \nhan{AN_TOAN}, lỗi lớn nhất là lẫn \nhan{QUAY_ROI} thành \nhan{DE_DOA} (\so{C0.v2.cm.qr.dd} câu).
Hai nhãn này giáp nhau thật: câu bám riết thường kèm một lời ép (``không gửi ảnh thì\dots''), nên ranh giới giữa
``quấy rối'' và ``dọa'' khó cả với người gán.

\subsection{Báo nhầm câu an toàn}

Ở C0, \so{C0.v2.fp.tong} trên \so{C0.at.n} câu \nhan{AN_TOAN} bị báo nguy hại: \so{C0.v2.cm.at.dd} câu thành
\nhan{DE_DOA}, \so{C0.v2.cm.at.ld} câu thành \nhan{LUA_DAO}, \so{C0.v2.cm.at.qr} câu thành \nhan{QUAY_ROI}. Xem các câu
này, mô hình thường bám vào \emph{từ khóa của tình huống nguy hại} mà bỏ qua ý của cả câu: chữ ``OTP'', ``phí ship'',
``báo cô'', ``chặn đường''. Bảng~\ref{tab:p2-bao-nham} là hai ví dụ.

\begin{table}[H]
\centering
\small
\caption{Hai câu an toàn bị v2 báo nhầm (C0).}
\label{tab:p2-bao-nham}
\begin{tabularx}{\linewidth}{>{\hsize=1.45\hsize}Y >{\hsize=0.55\hsize}Y}
\toprule
Tin nhắn & Kết quả \\
\midrule
\so{vd.otp.text} &
Đúng: \so{vd.otp.gold}\newline v2: \so{vd.otp.pred} (\so{vd.otp.tincay})\newline gpt-oss-120b: \so{vd.otp.llm}\newline v2 tô: \so{vd.otp.bc} \\
\midrule
\so{vd.baoco.text} &
Đúng: \so{vd.baoco.gold}\newline v2: \so{vd.baoco.pred} (\so{vd.baoco.tincay})\newline gpt-oss-120b: \so{vd.baoco.llm}\newline v2 tô: \so{vd.baoco.bc} \\
\bottomrule
\end{tabularx}
\end{table}

Câu đầu là tin OTP thật: nó không đòi gì, chỉ bảo nhập mã vào ứng dụng. Hướng dẫn gán nhãn ghi rõ loại tin này là
\nhan{AN_TOAN}, nhưng trong train, chữ ``OTP'' gần như chỉ xuất hiện ở đoạn lừa đảo. Câu sau là bạn bè khuyên nhau
\emph{báo cô giáo}, tức là cách xử lý đúng khi bị bắt nạt; mô hình lại thấy ``gửi cho cô'', ``gặp nhóm kia'' và đoán
dọa. Cả hai đều bị báo với độ tin cậy cao, nên chỉ hạ ngưỡng chung không đủ sửa.

\textbf{Tiếng lóng.} Câu tiếng lóng dễ bị báo nhầm vì chứa từ tục mà không nhắm ai. Trên dev, ở lát câu có thẻ tiếng
lóng (n = \so{dev.tieng.n}), mô hình chốt có FPR \so{dev.mmbert.tieng.fpr} ở lát này, ngang BamiBERT
(\so{dev.bamibert.tieng.fpr}), thấp hơn GuardMerge (\so{dev.guard.tieng.fpr}) nhưng cao hơn ViSoBERT
(\so{dev.visobert.tieng.fpr}). Bù lại, Macro-F1 của nó trên lát này cao nhất trong 4 ứng viên (\so{dev.mmbert.tieng} so với
\so{dev.visobert.tieng} của ViSoBERT), nên tổng thể vẫn là một lý do chọn mmBERT-base. Dù vậy, mức báo nhầm này vẫn cao.

\subsection{Bỏ sót câu nguy hại}

Ở C0, \so{C0.v2.fn.tong} câu nguy hại bị đoán \nhan{AN_TOAN}: \so{C0.v2.cm.dd.at} câu \nhan{DE_DOA}, \so{C0.v2.cm.qr.at}
câu \nhan{QUAY_ROI}, \so{C0.v2.cm.ld.at} câu \nhan{LUA_DAO} và \so{C0.v2.cm.xp.at} câu \nhan{XUC_PHAM}. Hơn nửa là câu
mang tính ẩn ý: lời dọa nói tránh (``mai đi đường đông người nhé''), kiểm soát người yêu bằng giọng quan tâm. Phần lớn
thuộc kịch bản chưa có trong train; ví dụ cụ thể ở phần Yêu cầu 2.

\subsection{Hiệu chỉnh ngưỡng: đã làm gì, chưa làm gì}

\begin{itemize}
  \item \textbf{Đã làm:} dò nhiệt độ $T$ và hai ngưỡng bằng chứng trên tập cal, khóa trước khi chạy test
        (Mục~\ref{sec:p2-mo-hinh}).
  \item \textbf{Chưa làm:} dò ngưỡng riêng cho từng nhãn (ví dụ đòi xác suất \nhan{DE_DOA} cao hơn mới báo). Đây là cách
        trực tiếp nhất để hạ FPR, và là việc tiếp theo. Cần một tập validation do người gán; dò trên cal (nhãn LLM) có
        nguy cơ khớp theo lỗi của LLM.
\end{itemize}

Kết quả thực nghiệm: FPR của v2 ở C0 (\so{C0.v2.fpr}) cao hơn gpt-oss-120b (\so{C0.llm.fpr}). Theo vùng, FPR cao nhất ở giọng Nam
(\so{C0.nam.v2.fpr}), so với Bắc \so{C0.bac.v2.fpr} và Trung \so{C0.trung.v2.fpr}. Mỗi vùng chỉ có khoảng
\so{C0.nam.v2.at.n}--\so{C0.bac.v2.at.n} câu an toàn, nên một câu sai làm FPR của vùng đổi vài điểm phần trăm; chênh lệch
vùng cần đọc kèm cỡ mẫu này.

\subsection{Chất lượng bằng chứng}

\begin{figure}[H]
\centering
% Sinh tự động bằng `python -m src.v2.eval.xuat_so_lieu_bao_cao`. Không sửa tay.
\begin{tikzpicture}
\begin{axis}[width=0.75\linewidth, height=5.8cm, ymin=0, ymax=100,
  symbolic x coords={C0,C1,C2}, xtick=data, enlarge x limits=0.15, ylabel={F1 bằng chứng (theo ý)},
  axis x line*=bottom, axis y line*=left, ymajorgrids, grid style={draw=black!12},
  tick label style={font=\small}, label style={font=\small}, clip=false,
  xticklabels={C0 (gốc),C1 (biến dạng thường),C2 (điểm mù)}]
\addplot[mauvhai, line width=2pt, mark=*, mark size=3pt] coordinates {(C0,83.0) (C1,76.5) (C2,75.3)};
\addplot[maullm, line width=2pt, mark=*, mark size=3pt] coordinates {(C0,70.7) (C1,56.1) (C2,32.2)};
\node[anchor=west,font=\small] at (axis cs:C2,75.3) {\ SafeViChi v2: 75,3};
\node[anchor=west,font=\small] at (axis cs:C2,32.2) {\ gpt-oss-120b: 32,2};
\node[anchor=south,font=\small] at (axis cs:C0,83.0) {83,0};
\node[anchor=north,font=\small] at (axis cs:C0,70.7) {70,7};
\end{axis}
\end{tikzpicture}

\caption{F1 bằng chứng (theo ý) khi tin nhắn bị biến dạng dần. gpt-oss-120b vẫn đoán được nhãn nhưng không còn chỉ
đúng chỗ; v2 cắt cụm ngay trên văn bản gốc nên giữ được.}
\label{fig:p2-bc-bien-dang}
\end{figure}

\begin{table}[H]
\centering
\small
\caption{Bằng chứng tô sáng của v2 và gpt-oss-120b. n BC: số câu nguy hại được chấm bằng chứng. Tô oan: câu
\nhan{AN_TOAN} bị tô ít nhất một cụm. Rỗng: câu nguy hại được báo đúng là nguy hại nhưng không có cụm nào. Sai chữ:
cụm không có nguyên văn trong tin nhắn.}
\label{tab:p2-bang-chung}
\footnotesize
\begin{tabular}{llccccccc}
\toprule
Tập & Hệ & n BC & BC P & BC R & BC F1 & Tô oan & Rỗng & Sai chữ \\
\midrule
\multirow{2}{*}{C0} & gpt-oss-120b & \so{C0.llm.nbc} & \so{C0.llm.bcp} & \so{C0.llm.bcr} & \so{C0.llm.bcf} & \so{C0.llm.toan} & \so{C0.llm.rong} & \so{C0.llm.khongnguyenvan} \\
 & v2 & \so{C0.v2.nbc} & \so{C0.v2.bcp} & \so{C0.v2.bcr} & \so{C0.v2.bcf} & \so{C0.v2.toan} & \so{C0.v2.rong} & \so{C0.v2.khongnguyenvan} \\
\midrule
\multirow{2}{*}{C1} & gpt-oss-120b & \so{C1.llm.nbc} & \so{C1.llm.bcp} & \so{C1.llm.bcr} & \so{C1.llm.bcf} & \so{C1.llm.toan} & \so{C1.llm.rong} & \so{C1.llm.khongnguyenvan} \\
 & v2 & \so{C1.v2.nbc} & \so{C1.v2.bcp} & \so{C1.v2.bcr} & \so{C1.v2.bcf} & \so{C1.v2.toan} & \so{C1.v2.rong} & \so{C1.v2.khongnguyenvan} \\
\midrule
\multirow{2}{*}{C2} & gpt-oss-120b & \so{C2.llm.nbc} & \so{C2.llm.bcp} & \so{C2.llm.bcr} & \so{C2.llm.bcf} & \so{C2.llm.toan} & \so{C2.llm.rong} & \so{C2.llm.khongnguyenvan} \\
 & v2 & \so{C2.v2.nbc} & \so{C2.v2.bcp} & \so{C2.v2.bcr} & \so{C2.v2.bcf} & \so{C2.v2.toan} & \so{C2.v2.rong} & \so{C2.v2.khongnguyenvan} \\
\bottomrule
\end{tabular}
\end{table}

\begin{itemize}
  \item Trên C0, bằng chứng của v2 vừa đúng hơn vừa đủ hơn gpt-oss-120b: P/R/F1 theo ý \so{C0.v2.bcp}/\so{C0.v2.bcr}/%
        \so{C0.v2.bcf} so với \so{C0.llm.bcp}/\so{C0.llm.bcr}/\so{C0.llm.bcf}.
  \item Khi chữ bị biến dạng, gpt-oss-120b gần như không trả được bằng chứng dùng được: ở C2, BC F1 chỉ còn
        \so{C2.llm.bcf} và \so{C2.llm.khongnguyenvan} số cụm không có nguyên văn trong tin. Lý do: LLM đọc hiểu được chữ
        lạ nhưng khi chép lại thì ``sửa'' chúng thành chữ thường, nên cụm không còn khớp tin gốc và không tô được.
        v2 cắt cụm từ chính văn bản nên vẫn giữ \so{C2.v2.bcf}.
  \item Điểm yếu của v2: \so{C0.v2.toan} câu an toàn bị tô oan (cao hơn \so{C0.llm.toan} của gpt-oss-120b). Con số này
        đi theo FPR: hễ câu an toàn bị báo nguy hại thì gần như luôn có cụm được tô.
\end{itemize}

\ghinho{\textbf{Cách chấm bằng chứng.} Điểm bằng chứng chỉ tính trên câu nguy hại mà hệ đã báo là nguy hại; câu nguy hại
bị đoán \nhan{AN_TOAN} đã bị phạt ở phần phân loại. Vì v2 có recall cao hơn, nó được chấm bằng chứng trên nhiều câu hơn
(C0: \so{C0.v2.nbc} so với \so{C0.llm.nbc} câu của gpt-oss-120b), tức là trên cả những câu khó mà LLM bỏ qua.}
